\documentclass[sigconf,nonacm,11pt]{acmart}

\usepackage{booktabs}
\usepackage{amsmath}
\usepackage{graphicx}

% colonnes etroites + identifiants longs non coupables
\setlength{\emergencystretch}{3em}

\settopmatter{printacmref=false}
\renewcommand\footnotetextcopyrightpermission[1]{}
\pagestyle{plain}

\begin{document}

\title{Bytes Are Not Milliseconds}
\subtitle{A Measurement Study of Two-Pass Block Selection for Long-Context LLM Inference}

\author{Anonymous}
\affiliation{\institution{}\country{}}

\begin{abstract}
Block-sparse attention makes long-context decoding tractable by reading only a
fraction of the KV cache, but the \emph{selector}---the mechanism that decides
which blocks to read---has itself become the dominant cost. We study
\emph{Adaptive Sequential Precision} (ASP), a two-pass selector that scores all
blocks with a coarse per-block summary, then rescores only the survivors with a
finer one, and we formalise it as a \emph{precision-allocation bandit}: a
best-arm-identification problem in which the measurement budget is spent in
\emph{bits per arm} along a rate--distortion curve rather than in repeated
stochastic pulls.

Our headline result is negative and, we argue, instructive. On Qwen3-8B, ASP
reduces selector KV traffic by $13\text{--}25\times$ yet is $1.2\text{--}1.7\times$
\emph{slower} end-to-end. We explain the gap in three measured steps. First, the
irregular gather that ASP requires is essentially free on datacenter hardware:
across three memory architectures (LPDDR5, GDDR6, HBM2e) scattered reads reach
$98\text{--}104\%$ of contiguous bandwidth once each gathered chunk exceeds
$\approx$1\,KiB, and on A100 the granularity threshold disappears entirely.
Second, fusing the two passes into a single kernel with interleaved stratified
selection yields $1.57\times$ over the two-kernel form on an H200 (15/15
configurations), but only when the per-stratum quota is small ($q\!\approx\!8$);
at $q=64$ fusion loses by $3\times$. Third, the residual end-to-end gap is
per-layer kernel-launch overhead: an unfused PyTorch selector costs a flat
$\approx$20\,ms per decode step, which alone exceeds the entire dense attention
time ($2\text{--}11$\,ms at $16$k--$262$k tokens, itself running at $84\%$ of
peak bandwidth).

We additionally report two \emph{refuted} predictive models of the gather
granularity threshold---one based on Little's law, one on DRAM page
invariance---and a measurement-driven law relating selection recall to
estimator noise, whose residuals are explained by Jensen's inequality over
per-query noise dispersion ($r=0.96$). We argue that convergent empirical
behaviour across three memory architectures without a reliable predictive model is itself a
result worth documenting.
\end{abstract}

\keywords{LLM inference, sparse attention, KV cache, GPU kernels, memory bandwidth}

\maketitle

%======================================================================
\section{Introduction}

Decoding a token from a long-context language model is bound by memory
bandwidth, not by arithmetic. Every generated token requires reading the entire
key--value (KV) cache, whose size grows linearly with context. At one million
tokens a mid-sized model must stream tens of gigabytes \emph{per token per
sequence}, and no amount of additional compute helps.

Block-sparse attention answers this by reading only the $k$ most relevant blocks
of the cache~\cite{yuan2025nsa,tang2024quest}. This shifts the problem rather
than removing it: to know which $k$ blocks matter, something must look at all
$n$ of them. That \emph{selector} is dense by construction, and it has quietly
become the dominant term. Reconstructing the cost of DeepSeek-V3.2 from its
published hyperparameters, we find that at 1M context its lightning indexer
accounts for $\approx$90\% of decode FLOPs and must stream $7.9$\,GB per decoded
token---while the sparse attention it enables costs $3\%$.\footnote{Derived
independently from the architecture description
in~\cite{deepseek2025v32,deepseek2024v3}, whose MLA latent cache originates
in~\cite{deepseek2024v2}; our reconstruction reproduces the
published parameter counts of DeepSeek-V4 to within $1.7\%$, which we take as
evidence the cost model is faithful.}

\paragraph{Why the existing approach is suboptimal.} Every selector we are aware
of---NSA~\cite{yuan2025nsa}, Quest~\cite{tang2024quest},
COBS~\cite{tian2026cobs}, DeepSeek's CSA~\cite{deepseek2026v4},
AsyncTLS~\cite{hu2026asynctls}---allocates the \emph{same} number of summary
bytes to every block and reads each summary exactly once. This is a uniform,
single-pass design. It is an assumption, not a necessity: most blocks are far
from the decision boundary and could be discarded from a much coarser summary.

\paragraph{Contribution.} We formalise non-uniform, sequential summary reading
as a \emph{precision-allocation bandit}, implement it as a single fused GPU
kernel with stratified selection, and measure it end-to-end on a real model---%
finding that the byte savings are real, large, and \emph{do not become
speedups} unless the selector is fused, and unless the surrounding engine is
itself memory-bound.

%======================================================================
\section{Background and Related Work}

\paragraph{Sparse attention and block selection.} Fixed-pattern methods attend
to predetermined windows and global tokens~\cite{beltagy2020longformer,%
zaheer2020bigbird}. Eviction methods discard low-importance
entries~\cite{xiao2024streamingllm,zhang2023h2o,li2024snapkv}. Query-aware block
selectors summarise contiguous blocks offline and, per query, run fine-grained
attention over the highest-scoring ones; this family includes
Quest~\cite{tang2024quest} and NSA~\cite{yuan2025nsa}. COBS~\cite{tian2026cobs}
formalises selection as ranking blocks by \emph{attention mass} and shows that
first-order (mean-pool) summaries are the limiting factor, proposing a
second-order cumulant correction. Our work is complementary and, on one point,
in tension: we find that at matched byte budget a \emph{support}-based summary
(a coreset of $r$ centroids scored by log-sum-exp) dominates a \emph{moment}-based
one (mean plus rank-$r$ covariance) at equal or lower cost: a coreset of $r=2$
centroids ($2.03$ vectors per block) reaches $86.4\%$ mass recall at $k=8$, where
COBS at $r=8$ ($9.12$ vectors) reaches $81.5\%$. The reason is that the
within-block key covariance is not low-rank---its effective rank is $17.8$ out of $64$ in our measurements, so
a rank-$r$ summary captures only a varying fraction of the quadratic term. The
same measurement bounds what methods built on low-rank key structure, such as
Loki~\cite{singhania2024loki}, can recover \emph{within} a block---though Loki
estimates its rotated basis globally, where the spectrum is more favourable.

\paragraph{Hierarchical selection.} AsyncTLS~\cite{hu2026asynctls} performs
coarse block filtering followed by fine token selection;
HiSparse~\cite{hisparse2026} hierarchises \emph{storage} between host and device,
and lookahead sparse attention~\cite{lsa2026} amortises the index across steps.
Both refine \emph{what} is examined. ASP is orthogonal: it refines \emph{how
precisely} the same objects are examined, and allocates that precision according
to distance from the decision threshold.

\paragraph{Serving systems.} PagedAttention~\cite{kwon2023vllm} and
RadixAttention~\cite{zheng2024sglang} govern KV memory;
FlashAttention-3~\cite{shah2024flashattention3} and
FlashInfer~\cite{ye2025flashinfer} provide the kernels; prefill/decode
disaggregation~\cite{zhong2024distserve,qin2024mooncake,agrawal2024sarathi}
governs placement. ASP sits below all of these, inside the attention kernel.

\paragraph{Speculative decoding.} EAGLE-3~\cite{li2025eagle3} and its
predecessors~\cite{leviathan2023speculative} amortise the KV read across
$\gamma+1$ candidate tokens. We note in passing---and quantify in
Section~\ref{sec:disc}---that sparse attention and speculation interact
negatively, because candidate positions select different blocks and the
verification pass must read their union.

\paragraph{Bandits.} Selecting the top-$k$ of $n$ noisy quantities under a
measurement budget is best-arm identification~\cite{bubeck2009pure,%
carpentier2016tight,karnin2013almost}. We import that framework, with one
structural difference developed in Section~\ref{sec:method}.

%======================================================================
\section{Method}\label{sec:method}

\subsection{Block summarisation as measure quantisation}

A block $b$ is an empirical measure on key space, $\mu_b=\sum_{j\in b}\delta_{k_j}$,
and its attention mass is the Laplace transform of that measure,
$m_b(q)=\int e^{s\,q^\top k}\,d\mu_b(k)$ with $s=1/\sqrt{D}$. Summarising a block
means replacing $\mu_b$ by a few-atom measure whose transform agrees on the query
distribution. Writing $\ell_b=\log m_b$, the Gibbs variational identity gives
\begin{equation}
\ell_b(q)=\max_{p\in\Delta_L}\Big[\textstyle\sum_r p_r\,(s\,q^\top k_r)+H(p)\Big],
\end{equation}
so $\ell_b$ interpolates between $\log L+\text{mean}$ as $q\to 0$ and
$\max_r s\,q^\top k_r$ as $\|q\|\to\infty$. Two approximation families follow:
\emph{moment} approximations expand the cumulant generating function about $q=0$
(mean-pool at order one, COBS at order two), while \emph{support} approximations
approximate the geometry of the key cloud.

\begin{proposition}[Uniform two-sided coreset bound]\label{prop:bound}
Partition a block into clusters $C_1,\dots,C_r$ with centroids $\mu_c$, counts
$n_c$ and radii $\rho_c=\max_{k\in C_c}\|k-\mu_c\|$. Then for \emph{every} query $q$,
\begin{equation}
\Big|\,\ell_b(q)-\log\textstyle\sum_c n_c e^{s\,q^\top\mu_c}\Big|\;\le\;s\,\|q\|\max_c \rho_c .
\end{equation}
\end{proposition}
\begin{proof}
For $k\in C_c$, Cauchy--Schwarz gives $|q^\top(k-\mu_c)|\le\|q\|\rho_c$, hence
$n_c e^{s q^\top\mu_c-s\|q\|\rho_c}\le\sum_{k\in C_c}e^{s q^\top k}\le
 n_c e^{s q^\top\mu_c+s\|q\|\rho_c}$. Summing over $c$ and taking logarithms, the
common factors $e^{\pm s\|q\|\rho}$ factor out.
\end{proof}

Unlike a truncated cumulant expansion, this bound is uniform in the direction of
$q$, is sign-aware (the term $e^{q^\top\mu_c}$ distinguishes $+\mu$ from $-\mu$,
whereas $q^\top\Sigma q$ does not), and is exact on an isolated outlier, which
forms a singleton of zero radius. By Jensen's inequality the coreset estimate is
moreover always a \emph{lower} bound on $\ell_b$, so the interval is one-sided
and half as wide as the proposition suggests. The bound also identifies the right
clustering objective: minimising the \emph{maximum} radius, i.e.\
$k$-center~\cite{gonzalez1985clustering}, rather than $k$-means.

\subsection{The precision-allocation bandit}

Let $n$ blocks be candidates and $k$ be selected. A summary of $b$ bytes yields an
estimator with noise $\sigma(b)$, a measurable rate--distortion curve---the same
curve that governs KV-cache quantization~\cite{turboquant2025,blockgtq2026}, applied
there uniformly to every block. A two-pass
scheme reads $b_1$ bytes for all $n$ blocks and $b_2$ bytes for $m$ survivors:
\begin{equation}
B(b_1,b_2,m)=n\,b_1+m\,b_2 .
\end{equation}
This is best-arm identification with a crucial difference: in a classical bandit,
pulling an arm $t$ times reduces noise as $1/\sqrt t$ by stochastic averaging;
here, reading $b$ bytes reduces it along a \emph{deterministic} rate--distortion
curve. We call this a \emph{precision-allocation bandit}. To our knowledge this
variant is treated neither in the bandit literature nor in the attention literature.

Recall loss decomposes into two first-order-independent sources. A block of true
rank $j\le k$ is wrongly eliminated in pass~1 if its noisy score falls below that
of the rank-$m$ block; both scores being independently perturbed, their difference
has standard deviation $\sigma_1\sqrt2$, giving
\begin{equation}
L_1(b_1,m)=\sum_{j\le k} w_j\,\Phi\!\Big(\!-\Delta_{j,m}\big/(\sigma(b_1)\sqrt2)\Big),
\label{eq:l1}
\end{equation}
with $\Delta_{j,m}=\ell_{(j)}-\ell_{(m)}$ and $w_j$ the mass share of rank $j$.
Pass-2 ranking error is $L_2(b_2)=1-R(\sigma(b_2))$, where $R$ is the empirical
recall-versus-noise curve of Section~\ref{sec:law}. The design problem is
\begin{equation}
\min_{b_1,b_2,m} \; n\,b_1+m\,b_2
\quad\text{s.t.}\quad R(\sigma(b_2))-L_1(b_1,m)\ \ge\ \tau .
\end{equation}
All quantities---$\sigma(b)$, $\Delta_{j,m}$, $w_j$, $R$---are measurable offline
on a calibration sample, so the configuration is computed, not swept.

\subsection{Fusion by interleaved stratification}\label{sec:fusion}

\begin{figure}[t]
\centering
\includegraphics[width=\columnwidth]{figs/strat.pdf}
\caption{Cost of stratified selection relative to the global top-$m$, versus the
number of strata $G$. Interleaved stratification (solid) is nearly free;
contiguous stratification (dashed) is not, because relevance is correlated
between adjacent blocks.}
\label{fig:strat}
\end{figure}


A naive fused kernel must satisfy a global dependency: pass~2 needs the top-$m$
of pass~1 across all blocks of a query. Three designs are possible. One work-group
per query needs no synchronisation but collapses parallelism to the number of
queries (batch $\times$ KV heads), typically two orders of magnitude below GPU
residency. A grid-wide barrier preserves parallelism but costs $5$--$10\,\mu$s on
A100---the same order as the launch it was meant to save.

We take a third route. ASP is a heuristic filter, not an exact ranking, so the
\emph{global} top-$m$ is not required. We partition the $n$ blocks into $G$
interleaved strata; work-group $g$ of a query handles blocks
$\{g,\,g\!+\!G,\,g\!+\!2G,\dots\}$, keeps the top-$q$ of its own stratum with
$q=m/G$, and immediately gathers its own survivors. There is \emph{no
synchronisation at all}, parallelism is complete, and there is a single launch.

Interleaving is not a detail. Measured on real attention maps (Figure~\ref{fig:strat}),
interleaved stratification costs at most $2.9$--$3.2$ relative recall points across
all $(m,G)$, whereas \emph{contiguous} stratification costs up to $22.3$ points:
adjacent blocks are correlated in relevance, so a contiguous stratum can hold a
disproportionate share of the mass and be truncated by its quota. Interleaving
breaks that correlation, and it is also the natural grid-stride access pattern.

%======================================================================
\section{Experimental Setup}

\paragraph{Hardware.} Four GPUs spanning three memory technologies: an AMD
Radeon 860M integrated GPU (\texttt{gfx1152}, RDNA 3.5, shared LPDDR5,
80\,GB/s measured), an NVIDIA Tesla T4 (GDDR6, 274\,GB/s measured), an NVIDIA
A100 PCIe 40\,GB (HBM2e, 1372\,GB/s measured) and an NVIDIA H200 NVL
(143\,GB, 132 SMs). Kernels are OpenCL on the AMD part and CUDA
(\texttt{nvcc -O3}) elsewhere; a single source compiles under both \texttt{nvcc}
and \texttt{hipcc}, using no warp-width-dependent primitive.

\paragraph{Selection-quality measurements.} Real $Q$/$K$ activations captured from
Qwen2.5-0.5B and SmolLM2-135M on WikiText at 8192 tokens. Following the
convention of block-sparse systems, the local window and attention sink are
served by separate branches and are therefore \emph{excluded} from the candidate
set; without this exclusion the metric is dominated by trivially-retrieved blocks.
The metric is \emph{mass recall}: the true attention mass captured by the $k$
selected blocks, divided by that of the oracle top-$k$.

\paragraph{End-to-end.} Qwen3-8B~\cite{qwen3} (36 layers, 32 query heads, 8 grouped-query
KV heads~\cite{ainslie2023gqa}, head dim 128, 144\,KiB/token in bf16) on the H200, with ASP registered
through the \texttt{ALL\_ATTENTION\_\allowbreak FUNCTIONS} extension point of
\texttt{transformers}~5.16---the supported hook, not a monkey-patch. The rest of
the model runs unmodified, so the measurement is a \emph{complete decode step}.

%======================================================================
\section{Results}

\subsection{A law relating recall to estimator noise}\label{sec:law}

\begin{figure}[t]
\centering
\includegraphics[width=\columnwidth]{figs/law.pdf}
\caption{Reference curve $R(\sigma)$, obtained by injecting controlled Gaussian
noise into the \emph{exact} block log-masses and re-ranking. Three real summaries
are placed at their measured $\sigma$: coreset and mean-pool land on the curve,
COBS beats it by 5--8 points.}
\label{fig:law}
\end{figure}


We first establish what any selector can achieve. Injecting controlled Gaussian
noise of standard deviation $\sigma$ into the \emph{exact} block log-masses and
ranking gives a reference curve $R(\sigma)$ (Figure~\ref{fig:law}) that is
independent of any method: $95\%$ mass recall at $k=8$ requires
$\sigma\le0.39$~nats, and $99\%$ requires $\sigma\le0.16$.

Real methods do not all fall on this curve. Coresets and mean-pool fall within $0.9$
points of the prediction; COBS and Gaussian-mixture summaries beat it by $4.8$ to
$8.1$ points. We tested and rejected two explanations---a rank-dependent bias
(injecting the measured bias makes predictions \emph{worse}, because the
rank-conditional mean of the error is a regression-to-the-mean artefact) and
over-counting of between-stratum variance (the effect is negligible). The
explanation that survives is Jensen's inequality over \emph{per-query} noise
dispersion: $R$ is convex, so $\mathbb{E}_i R(\sigma_i)>R(\mathbb{E}_i\sigma_i)$,
and the gap correlates with the coefficient of variation of $\sigma$ across
queries at $r=+0.956$. The law is therefore correct but must be applied per query
and then averaged.

\paragraph{Blocks are nearly tied.} The median log-mass gap between the 8th and
9th best block is $0.070$~nats, against a residual estimator uncertainty of
$1.08$~nats without RoPE and $2.36$ with---a ratio of $15$ to $34$. No compact summary can resolve the ordering at
the cut. This single fact explains four independent negative results we obtained
(quadrature-optimised weights, correctness certificates, learned calibration and
adaptive $k$ all reduce estimation error without improving ranking), and it
explains why mass recall ($85$--$90\%$) so far exceeds set recall ($60$--$75\%$):
mis-selecting a near-tied block is nearly free.

\subsection{The irregular gather is free}

\begin{figure*}[t]
\centering
\includegraphics[width=\textwidth]{figs/gather.pdf}
\caption{Scattered versus contiguous reads at equal volume, on three memory
architectures. \textbf{Left:} bandwidth ratio. \textbf{Right:} fraction of peak
bandwidth reached by the contiguous path; below $\approx$30\% neither path
saturates and the ratio is uninformative. The gather penalty vanishes above
$\approx$1\,KiB and decreases monotonically with GPU parallelism; on A100 the
threshold has disappeared.}
\label{fig:gather}
\end{figure*}

ASP's second pass reads a query-dependent subset of summaries, an irregular
gather. The concern is misposed: each gathered element is a whole block summary
of $r\!\cdot\! D\!\cdot\!2$ bytes---1\,KiB for $r=8$, $D=64$, i.e.\ sixteen
64-byte cache lines---so coalescing is perfect \emph{within} a chunk and only
inter-chunk locality is lost.

Figure~\ref{fig:gather} confirms this and locates the threshold. In the saturated
regime (contiguous path above $80\%$ of peak) the gather averages $104.0\%$ of
contiguous bandwidth on the integrated GPU, $98.4\%$ on T4 and $97.7\%$ on A100,
and never falls below $94.9\%$ on any of the three; values above unity are
measurement noise on a path that is already bandwidth-saturated. Below saturation
the three separate: the integrated GPU falls to $67.0\%$ at 64-byte chunks and T4
to $88.0\%$ at 512 bytes, whereas A100 never drops below $95.5\%$ at \emph{any}
granularity---the threshold has disappeared. The design rule is
$r_2\!\cdot\! D\!\cdot\!2\ge2048$ bytes. The coarse pass violates it by
construction ($r_1=2$, $D=64$: 256 bytes), which is harmless precisely because
that pass reads every block and is therefore a contiguous scan, not a gather.

\paragraph{Two refuted models.} We pre-registered a predictive model of the
threshold based on Little's law~\cite{little2008queueing}: a gather kernel
saturates when $N_{\text{res}}\cdot\min(C,U)\ge B\lambda$, giving
$C^\star=\min(U,\alpha B\lambda/N_{\text{res}})$ with $\alpha$ calibrated on the
integrated GPU. It predicted $328$~bytes for T4---\emph{lower} than the 1024
observed there---and $1009$ for A100. Measurement refuted the calibration twice:
T4's threshold is $1024$, and A100 has none. We then adopted a DRAM-page
invariance hypothesis (the threshold is $\approx$1\,KiB on any memory system);
A100 refutes that too. What survives is only the \emph{direction} of the first
model: the penalty decreases monotonically with parallelism ($0.67$ at 64\,B on
8 CUs, $0.88$ at 512\,B on 40 SMs, $1.00$ at 512\,B on 108 SMs). We report this
as a result: the empirical behaviour converges across three memory architectures without
a predictive model that survives contact with the data, and we prefer to say so
rather than fit a third model post hoc.

\subsection{Fusion works, at small quota}

\begin{figure}[t]
\centering
\includegraphics[width=\columnwidth]{figs/fusion.pdf}
\caption{H200, $Q=64$ queries, real grids. Fusion beats the two-kernel scheme
only when the per-stratum quota is small; labels give winning configurations.}
\label{fig:fusion}
\end{figure}

\begin{table}[t]
\caption{Fused versus two-kernel ASP on H200 (59 configurations, $Q=64$,
$G\in\{32,128,512\}$).}
\label{tab:fusion}
\begin{tabular}{rrrr}
\toprule
$q$ & fused / two-kernel & fused / dense & winning \\
\midrule
\textbf{8}  & \textbf{1.57$\times$} & \textbf{1.71$\times$} (max 4.58) & \textbf{15/15} \\
16 & 0.78$\times$ & 1.27$\times$ & 3/15 \\
32 & 0.68$\times$ & 0.71$\times$ & 0/15 \\
64 & 0.32$\times$ & 0.38$\times$ & 0/14 \\
\bottomrule
\end{tabular}
\end{table}

Table~\ref{tab:fusion} and Figure~\ref{fig:fusion} give the result. Since
$G\cdot q=m$ is fixed, $G$ governs parallelism and $q$ the serial work per team;
the optimum is reached when the per-team quota $q/N_T$ falls to one, i.e.\
$q\approx N_T=8$, giving the design rule $G\approx m/8$. At that point fusion
beats the two-kernel form in 15 of 15 configurations, by $1.57\times$ median.

Reaching this required two diagnoses. The first fused kernel failed with an
illegal memory access, isolated by bisection to partial initialisation of a
shared array whose full extent was later scanned. The second ran but was
$3\times$ slower, with a runtime \emph{exactly linear in $q$} and independent of
summary size---the signature of serialisation, traced to a final merge of
per-team candidate lists executed by a single thread. Removing that merge by
stratifying one level further (each team gathers its own top-$q/N_T$) moved the
result from $0.31\times$ to $0.96\times$, and tuning $G$ to $1.57\times$.

\subsection{End to end: byte savings do not become speedups}

\begin{figure*}[t]
\centering
\includegraphics[width=\textwidth]{figs/e2e.pdf}
\caption{Qwen3-8B on H200. \textbf{Left:} full decode-step latency and the
corresponding byte-traffic ratio. \textbf{Right:} ventilation---dense attention
runs at $84\%$ of peak bandwidth and costs $2\text{--}11$\,ms, while the unfused
ASP selector costs a flat $20\text{--}23$\,ms.}
\label{fig:e2e}
\end{figure*}

\begin{table}[t]
\caption{Qwen3-8B, H200, complete decode step.}
\label{tab:e2e}
\small
\setlength{\tabcolsep}{4.5pt}
\begin{tabular}{rrrrrr}
\toprule
context & KV (GB) & dense & ASP & speedup & bytes \\
        &         & (ms)  & (ms) &         & saved \\
\midrule
16\,k  &  2.2 & 29.9 & 49.6 & 0.60$\times$ & 13.5$\times$ \\
65\,k  &  9.0 & 29.8 & 49.5 & 0.60$\times$ & 20.9$\times$ \\
131\,k & 18.0 & 30.0 & 49.4 & 0.61$\times$ & 23.0$\times$ \\
262\,k & 36.0 & 42.1 & 54.0 & 0.78$\times$ & 24.2$\times$ \\
524\,k & 72.0 & 72.6 & 89.8 & 0.81$\times$ & 24.9$\times$ \\
\bottomrule
\end{tabular}
\end{table}

Table~\ref{tab:e2e} is the paper's central measurement. ASP reads $13$--$25\times$
fewer bytes and is $1.2$--$1.7\times$ slower. We had predicted a $10$--$40\%$ gain;
we measure a $19$--$40\%$ loss.

The ventilation (Figure~\ref{fig:e2e}, right) identifies the cause completely.
Dense attention costs $2.06$, $4.06$, $6.62$ and $11.34$\,ms at $16$k, $65$k,
$131$k and $262$k---about $84\%$ of the peak-bandwidth floor, i.e.\ already
near-optimal. The unfused ASP selector costs a flat $20$--$23$\,ms: eight
non-fused PyTorch operations per layer, some 290 kernel launches per decode step.
Crossover requires dense attention to exceed the selector overhead, which happens
at $N\approx470$k; at $524$k we still measure $0.81\times$.

The prediction failed because it assumed a memory-bound decode. This integration
is not: the measured step is $3.2\times$ above its memory floor ($30$\,ms against
$9.4$\,ms at $131$k), the non-attention portion being dominated by framework
overhead. In an engine with CUDA graphs and fused kernels that portion would fall
towards $\approx$5\,ms, attention would dominate, and the byte savings would
matter---provided the selector is fused. That projection is an inference, not a
measurement, and we mark it as such.

\paragraph{An artefact that would have inverted the conclusion.} Our first
ventilation reported a $10.3\times$ ASP speedup on attention. It was wrong: both
paths materialised the expanded KV via \texttt{repeat\_interleave} to emulate
grouped-query attention~\cite{ainslie2023gqa}, inflating the dense baseline by the group factor $g=4$
while barely touching ASP, which gathers a tiny KV. The impossibility was visible
---attention alone ($121$\,ms) exceeded the full step ($30$\,ms). Enabling
GQA-aware SDPA on both paths removed it. We report this because the artefact
favoured our own method, and because the check that caught it (a component
exceeding its container) is cheap and general.

%======================================================================
\section{Discussion}\label{sec:disc}

\paragraph{Where the effort should go.} Two measurements reorder priorities.
First, at $k=8$, moving from a coreset selector to a perfect
oracle recovers $5.4$ points of output error, whereas doubling $k$ recovers
$8.5$. Second, the near-tie structure caps what any selector can do. Both suggest
that making $k$ larger but cheaper is worth more than making the selector smarter.

\paragraph{Sparse attention and speculation conflict.} Under sparse attention the
$\gamma+1$ speculative candidates select different blocks, so verification must
read their union. Measured on real attention, amortisation efficiency falls to
$41.8\%$ at $\gamma=8$ with token-granular selection, and speculative speedup can
vanish entirely ($1.00\times$ at $\alpha=0.7$). Block granularity partially
rescues it ($65.2\%$), which is an argument for block selection that its
proponents do not make.

\paragraph{Should content-defined blocking be in this paper?} We also find that
partitioning the cache by key \emph{geometry} rather than by position doubles the
attention mass captured by the top-8 blocks ($21.6\%\to43.1\%$ without RoPE).
We deliberately exclude it here. It shares the theoretical frame---the radius
$\rho$ of Proposition~\ref{prop:bound} bounds every summary's error---but answers
a different question (\emph{how to partition}, not \emph{how to allocate
precision}), and its validation requires work this paper does not contain: the
gain is measured at a single $k$ and a single block size, and it collapses once
RoPE is applied---from $\times1.99$ to $\times1.23$ ($44.6\%\to54.8\%$), because
rotation partly destroys the geometric clustering the partition exploits.
Clustering pre-RoPE keys while attending post-RoPE recovers only part of it
($\times1.10$), a tension we have characterised but not resolved. Merging the two would dilute both.

%======================================================================
\section{Limitations}

\textbf{Scope of the hardware claims.} Four GPUs, one of them integrated. No
MI300 or Blackwell part was tested; the granularity threshold has no surviving
predictive model, so extrapolation is unwarranted.

\textbf{The end-to-end integration is not a serving engine.} We patch
\texttt{transformers}, not vLLM or SGLang: no continuous batching, no paged KV,
no CUDA graphs. Our central negative result is therefore conditional on an
integration that is itself $3.2\times$ from its memory floor. The fused kernel of
Section~\ref{sec:fusion} was measured \emph{outside} the model; joining the two is
the obvious next step and we have not done it.

\textbf{Context beyond the trained range.} Qwen3-8B has
$\texttt{max\_position\_embeddings}=40960$. Beyond $40$k we measure performance,
not quality; the memory pattern and latency are unchanged but outputs are not
meaningful. Selection quality was measured separately, on real attention, at
in-distribution lengths.

\textbf{Synthetic KV values and simplified summaries.} Decode latency depends on
cache \emph{shape}, not content, so we synthesise KV beyond $131$k; summaries in
the end-to-end path are segment means rather than $k$-center coresets, a weaker
but cheaper construction whose \emph{size}---all that latency depends on---is
identical.

\textbf{Quality of the fused selector.} Fusion's optimum sits at per-stratum
quota one, which costs $2.5$--$3.2$ relative recall points. We measured that cost
but did not measure its effect on any downstream benchmark such as
RULER~\cite{hsieh2024ruler}.

\textbf{Bandit assumptions.} Equation~\eqref{eq:l1} assumes independence between
the two loss sources and Gaussian score noise. The first is a first-order
approximation; the second is contradicted by our own finding that per-query noise
dispersion matters (Section~\ref{sec:law}).

%======================================================================
\section{Conclusion and Future Work}

Two-pass precision allocation reduces selector traffic by $13$--$25\times$ on a
real model, and that reduction does not become a speedup unless two conditions
hold: the selector must be a single fused kernel, and the surrounding engine must
be memory-bound. We measured the first condition being met ($1.57\times$ over the
two-kernel form on H200) and the second being violated ($3.2\times$ from the
memory floor). The irregular gather, which motivated this study, turns out to
cost nothing on datacenter hardware.

Four threads remain open. \emph{Fusing the validated kernel into a serving
engine} is the direct continuation and would test the projection of
Section~5.4. \emph{Content-defined blocking} needs its own study, with the
layer-dependence of its gain as the first question. \emph{The second error
mechanism}---the rank structure of estimator error, which a permutation test
shows is actively harmful by $1.2$--$5.2$ points---is characterised but
unexplained. And the \emph{precision-allocation bandit} deserves a proper
optimality theorem: we have a first-order design rule and a measured
rate--distortion curve, not a lower bound.

\bibliographystyle{ACM-Reference-Format}
\bibliography{refs}

\end{document}
